\section{The sum of all correlations reveals a square}
Let $r=\sum_{j=0}^{n-1}a_j$ be the row sum. Add every periodic correlation, including shift zero. By definition, each $C_s$ is itself a sum over $j$, so the total is a double sum. A finite sum may be added in any order, so we may add over $s$ first for each fixed $j$, and then pull out the factor $a_j$, which does not depend on $s$:
\begin{align*}
\sum_{s=0}^{n-1}C_s
&=\sum_{s=0}^{n-1}\sum_{j=0}^{n-1}a_j a_{(j+s)\bmod n}
 =\sum_{j=0}^{n-1}\sum_{s=0}^{n-1}a_j a_{(j+s)\bmod n}\\
&=\sum_{j=0}^{n-1}a_j\left(\sum_{s=0}^{n-1}a_{(j+s)\bmod n}\right)\\
&=\sum_{j=0}^{n-1}a_j\left(\sum_{k=0}^{n-1}a_k\right)\\
&=r^2.
\end{align*}
In the second line, for each fixed $j$, the indices $(j+s)\bmod n$ run through every position exactly once. This is the specific symmetry that justifies the reindexing.

One way to see the whole count is to use three abstract entries $a_0,a_1,a_2$. The three periodic correlations list these products:
\begin{align*}
 C_0&=a_0a_0+a_1a_1+a_2a_2,\\
 C_1&=a_0a_1+a_1a_2+a_2a_0,\\
 C_2&=a_0a_2+a_1a_0+a_2a_1.
\end{align*}
Together they contain each ordered product $a_j a_k$ once. Expanding $(a_0+a_1+a_2)^2$ gives those same nine ordered products, so the totals agree. For general $n$, each fixed first index $j$ meets every second index exactly once as the shift varies. That is why the total is $(\sum_j a_j)^2$, rather than, for example, $\sum_j a_j^2$.

This identity itself does not require signs; the same finite rearrangement holds for real entries. The restriction to signs is used afterward to identify $C_0=n$ and to make the row sum an integer. Separating those roles tells us which part of the argument might generalize.

If all nonzero $C_s$ vanish, the left side is just $C_0=n$. Hence
\[
n=r^2.
\tag{14.2}\label{eq:squareorder}
\]
The integer $r$ is a sum of signs, so the order $n$ must be a square. This is a proved necessary condition, valid for every positive length, not a pattern extrapolated from an experiment.
